\chapter{Varyans Analizi ve Grup Karşılaştırmaları}
\label{ch:anova}
\index{varyans analizi}
\index{ANOVA}

\begin{hedefler}
Bu bölümün sonunda okuyucu:
\begin{itemize}
  \item üç veya daha fazla grup için neden tek yönlü varyans analizine gereksinim duyulduğunu açıklayabilecek,
  \item ANOVA modelindeki grup içi ve gruplar arası değişkenliği ayırt edebilecek,
  \item kareler toplamı, serbestlik derecesi, ortalama kare ve $F$ istatistiğini hesaplayıp yorumlayabilecek,
  \item genel $F$ testi ile planlı ve post-hoc karşılaştırmaların farklı rollerini açıklayabilecek,
  \item klasik ANOVA ile Welch ANOVA arasında veri yapısına uygun seçim yapabilecek,
  \item eta-kare ve omega-kare etki büyüklüklerini hesaplayıp bağlam içinde yorumlayabilecek,
  \item bağımsızlık, normallik, varyans benzerliği ve aykırı değer koşullarını değerlendirebilecek,
  \item Python ve SPSS çıktılarından eksiksiz bir ANOVA bulgusu raporlayabilecektir.
\end{itemize}
\end{hedefler}

\begin{hazirlik}
Dört grubun bütün ikililerini karşılaştırmak için kaç ayrı $t$ testi gerekir?
Her test yüzde 5 anlamlılık düzeyinde yürütüldüğünde aile düzeyindeki yanlış
pozitif riskinin neden yüzde 5 olarak kalmadığını tartışın.
\end{hazirlik}

\section{PDR vakası: Üç destek programını karşılaştırmak}

Bir üniversite, sınav kaygısını azaltmak amacıyla üç destek yaklaşımını
karşılaştırmaktadır:

\begin{description}
  \item[A programı] Yapılandırılmış grup danışmanlığı,
  \item[B programı] Psikoeğitim semineri,
  \item[C programı] Çevrim içi öz-yardım materyali.
\end{description}

Her programa 12 öğrenci katılmıştır. Program sonunda düşük puanın daha az
kaygıyı gösterdiği bir ölçek uygulanmıştır. Grup ortalamaları ve standart
sapmaları şöyledir:

\begin{table}[htbp]
\centering
\caption{Üç programın son-test kaygı puanları.}
\begin{tabular}{@{}lccc@{}}
\toprule
Program & $n$ & Ortalama & Standart sapma \\
\midrule
A: Grup danışmanlığı & 12 & 44,83 & 4,61 \\
B: Psikoeğitim & 12 & 41,00 & 4,07 \\
C: Öz-yardım & 12 & 37,50 & 4,54 \\
\bottomrule
\end{tabular}
\end{table}

Örneklemde ortalamalar farklı görünmektedir. Ancak araştırma sorusu yalnız bu
36 öğrenciyi betimlemek değil, üç yaklaşımın karşılık geldiği evren
ortalamalarının farklı olup olmadığını değerlendirmektir. Tek yönlü ANOVA,
örneklem farklarını grup içindeki doğal değişkenlikle karşılaştırır.

\begin{researchbox}
Programlara öğrenciler kendi tercihiyle katılmışsa ortalama farklar programın
nedensel etkisini tek başına göstermez. Başlangıç kaygısı, zaman uygunluğu,
motivasyon veya yardım arama tutumu hem program seçimiyle hem sonuçla ilişkili
olabilir. Nedensel yorum için atama süreci kritik önemdedir.
\end{researchbox}

\section{Neden çok sayıda t testi yerine ANOVA?}
\index{çoklu karşılaştırma}
\index{aile düzeyinde hata oranı}

$k$ grubun bütün ikilileri için
\[
\frac{k(k-1)}{2}
\]
karşılaştırma gerekir. Üç grup için 3, dört grup için 6, beş grup için 10 ayrı
ikili test oluşur. Her birini $\alpha=0{,}05$ düzeyinde yürütmek, en az bir
yanlış pozitif sonuç bulma olasılığını artırır.

Testlerin bağımsız olduğu basit durumda üç test için aile düzeyindeki yanlış
pozitif olasılığı yaklaşık
\[
1-(1-0{,}05)^3\approx0{,}143
\]
yani yüzde 14,3'tür. Gerçek karşılaştırmalar ortak grupları kullandığı için
bağımsız değildir; yine de ayrı ayrı testlerin hata oranını artırdığı temel
sonuç değişmez.

Tek yönlü varyans analizi bütün grup ortalamalarının eşit olduğu ortak sıfır
hipotezini tek bir genel testle değerlendirir:
\[
H_0:\mu_1=\mu_2=\cdots=\mu_k.
\]
Alternatif hipotez ``bütün ortalamalar birbirinden farklıdır'' değildir; en az
bir evren ortalamasının diğerlerinden farklı olmasıdır.

\begin{mistakebox}
Anlamlı genel $F$ testi üç grubun da birbirinden farklı olduğunu göstermez.
Yalnız bütün ortalamaların eşit olduğu ortak hipotezle verinin uyumsuz
olduğunu söyler. Farkın hangi gruplar arasında bulunduğu ayrı karşılaştırmalarla
incelenir.
\end{mistakebox}

\section{Tek yönlü ANOVA modeli}
\index{varyans analizi!tek yönlü}
\index{gruplar arası değişkenlik}
\index{grup içi değişkenlik}

$j$ grubundaki $i$'inci kişinin puanı
\[
Y_{ij}=\mu+\tau_j+\varepsilon_{ij}
\]
biçiminde düşünülebilir. Burada $\mu$ genel ortalama, $\tau_j$ grubun genel
ortalamadan sapması ve $\varepsilon_{ij}$ kişinin kendi grup ortalaması
çevresindeki sapmasıdır.

Eşdeğer olarak her gözlemin genel ortalamadan sapması
\[
Y_{ij}-\bar Y_{..}
=\left(\bar Y_{.j}-\bar Y_{..}\right)
+\left(Y_{ij}-\bar Y_{.j}\right)
\]
şeklinde iki parçaya ayrılır:

\begin{itemize}
  \item grubun genel ortalamadan sapması: gruplar arası değişkenlik,
  \item kişinin kendi grup ortalamasından sapması: grup içi değişkenlik.
\end{itemize}

ANOVA'nın adı varyans analizidir; fakat temel araştırma sorusu ortalama
farklarıdır. Ortalamaların birbirinden ayrılması gruplar arası değişkenliği
büyüttüğü için soru varyans oranı üzerinden sınanır
\cite{maxwell2018,kutner2005}.

\begin{sezgibox}
Gruplar gerçekten aynı evrenden geliyorsa grup ortalamaları birbirine, her
gruptaki öğrenciler de kendi ortalamalarına benzer ölçüde dağılır. Grup
ortalamaları grup içindeki doğal saçılmaya göre çok uzaksa ortak ortalama
hipotezi inandırıcılığını kaybeder.
\end{sezgibox}

\section{Değişkenliğin ayrıştırılması}
\index{kareler toplamı}
\index{ortalama kare}
\index{F istatistiği@$F$ istatistiği}
\index{serbestlik derecesi}

Toplam kareler toplamı üç bileşenin temel ilişkisiyle ayrılır:
\[
SS_T=SS_A+SS_E.
\]

Burada
\[
SS_T=\sum_{j=1}^k\sum_{i=1}^{n_j}
(Y_{ij}-\bar Y_{..})^2
\]
toplam değişkenliği,
\[
SS_A=\sum_{j=1}^k n_j
(\bar Y_{.j}-\bar Y_{..})^2
\]
gruplar arası değişkenliği ve
\[
SS_E=\sum_{j=1}^k\sum_{i=1}^{n_j}
(Y_{ij}-\bar Y_{.j})^2
\]
grup içi ya da hata değişkenliğini gösterir.

Serbestlik dereceleri
\[
sd_A=k-1,
\qquad
sd_E=N-k,
\qquad
sd_T=N-1
\]
olup ortalama kareler
\[
MS_A=\frac{SS_A}{k-1},
\qquad
MS_E=\frac{SS_E}{N-k}
\]
şeklinde hesaplanır. Genel test istatistiği
\[
F=\frac{MS_A}{MS_E}
\]
oranıdır.

$H_0$ doğruysa gruplar arası ve grup içi ortalama kareler aynı hata
varyansını tahmin eder; $F$ çoğunlukla 1 çevresindedir. Grup ortalamaları
birbirinden uzaklaştıkça pay büyür ve $F$ değeri 1'in üzerine çıkar. $F$
negatif olamaz.

\section{Adım adım çözümlü ANOVA}

Üç program örneğinde toplam $N=36$, grup sayısı $k=3$ ve genel ortalama
$\bar Y_{..}=41{,}11$'dir. Ham puanlardan hesaplanan kareler toplamı:
\[
SS_A=322{,}89,
\qquad
SS_E=642{,}67,
\qquad
SS_T=965{,}56.
\]
Toplama denetimi $322{,}89+642{,}67=965{,}56$ sonucunu verir.

\subsection*{1. Serbestlik dereceleri}
\[
sd_A=3-1=2,
\qquad
sd_E=36-3=33,
\qquad
sd_T=36-1=35.
\]

\subsection*{2. Ortalama kareler}
\[
MS_A=\frac{322{,}89}{2}=161{,}44,
\qquad
MS_E=\frac{642{,}67}{33}=19{,}47.
\]

\subsection*{3. F istatistiği}
\[
F=\frac{161{,}44}{19{,}47}=8{,}29.
\]

\begin{table}[htbp]
\centering
\caption{Üç program için tek yönlü ANOVA tablosu.}
\begin{tabular}{@{}lrrrrr@{}}
\toprule
Kaynak & $SS$ & $sd$ & $MS$ & $F$ & $p$ \\
\midrule
Programlar arası & 322,89 & 2 & 161,44 & 8,29 & 0,001 \\
Grup içi & 642,67 & 33 & 19,47 & & \\
Toplam & 965,56 & 35 & & & \\
\bottomrule
\end{tabular}
\end{table}

Sonuç $F(2,33)=8{,}29$, $p=0{,}001$ olarak raporlanır. Bütün evren
ortalamalarının eşit olduğu hipoteze karşı kanıt vardır. Ancak henüz hangi
programların farklı olduğu bilinmemektedir.

\section{ANOVA varsayımları}
\index{varyans analizi!varsayımlar}
\index{bağımsızlık}
\index{normallik}
\index{varyans homojenliği}

Klasik tek yönlü ANOVA şu koşullara dayanır:

\begin{description}
  \item[Bağımsızlık] Bir gözlemin hata terimi diğer gözlemlerinkinden bağımsızdır.
  \item[Yaklaşık normallik] Her gruptaki hata veya artık dağılımı yaklaşık normaldir.
  \item[Varyans benzerliği] Grupların hata varyansları yeterince benzerdir.
  \item[Uygun ölçüm] Sonuç değişkeninde ortalama ve farklar anlamlıdır.
  \item[Uygun örnekleme] Seçim ve atama süreci hedeflenen yorumu destekler.
\end{description}

Bağımsızlık en kritik koşuldur ve bir normallik testiyle doğrulanamaz. Aynı
öğrencinin birden fazla ölçümü, öğrencilerin sınıflar içinde kümelenmesi veya
eşleştirilmiş katılımcılar farklı modeller gerektirir.

Normallik ham puanların bütün olarak değil, model artıklarının grup bazında
davranışına ilişkindir. Dengeli ve orta büyüklükte gruplarda ANOVA ılımlı
normallik sapmalarına dayanıklıdır. Küçük gruplarda ağır çarpıklık ve aykırı
değerler daha ciddi sorun oluşturur.

\begin{assumptionbox}
Önce her grubun $n$, ortalama, standart sapma, histogram, kutu grafiği ve
aykırı değerlerini inceleyin. Yalnız Shapiro--Wilk ve Levene testlerinin
$p$-değerleriyle yöntem seçmeyin; örneklem büyüklüğü ve sapmanın niteliğini de
değerlendirin.
\end{assumptionbox}

Şekil~\ref{fig:anova-yontem-secimi}, tek yönlü bağımsız grup karşılaştırmasında
tasarım ve veri yapısını birlikte kullanan yöntem seçimini özetler. Akış,
tek bir varsayım testinin $p$-değerini otomatik karar kuralı yapmaz.

\begin{figure}[H]
\centering
\begin{tikzpicture}[
  font=\small,
  decision/.style={diamond, draw=PrismBlue!80!black, fill=PrismBlue!10,
    align=center, aspect=2.35, inner sep=1.2pt, text width=3.25cm},
  result/.style={draw=PrismBlue!70!black, rounded corners=2pt,
    fill=PrismLight, align=center, minimum height=9mm, text width=3.1cm},
  arrow/.style={-{Stealth[length=2mm]}, thick, draw=PrismBlue!80!black},
  lab/.style={font=\footnotesize, fill=white, inner sep=1pt},
  node distance=7mm and 11mm
]
\node[decision] (design) {En az üç bağımsız grup ve sürekli sonuç mu?};
\node[result, right=of design] (other) {Eşleştirilmiş, tekrarlı veya kümeli tasarıma uygun model};
\node[decision, below=of design] (ordinal) {Sıralı ölçüm, ağır aykırı değer veya çok güçlü çarpıklık mı?};
\node[result, right=of ordinal] (robust) {Kruskal--Wallis ya da hedefi koruyan sağlam / yeniden örneklemeli yöntem};
\node[decision, below=of ordinal] (variance) {Varyanslar belirgin farklı veya grup büyüklükleri dengesiz mi?};
\node[result, right=of variance] (welch) {Welch ANOVA};
\node[result, below=of variance] (classic) {Klasik tek yönlü ANOVA};
\draw[arrow] (design) -- node[lab, above] {hayır} (other);
\draw[arrow] (design) -- node[lab, right] {evet} (ordinal);
\draw[arrow] (ordinal) -- node[lab, above] {evet} (robust);
\draw[arrow] (ordinal) -- node[lab, right] {hayır} (variance);
\draw[arrow] (variance) -- node[lab, above] {evet} (welch);
\draw[arrow] (variance) -- node[lab, right] {hayır} (classic);
\end{tikzpicture}
\caption{Tek yönlü grup karşılaştırmasında yöntem seçimi.}
\label{fig:anova-yontem-secimi}
\end{figure}

\section{Varyans eşitsizliği ve Welch ANOVA}
\index{Welch ANOVA}
\index{varyans eşitsizliği}
\index{Brown--Forsythe testi}
\index{Levene testi}

Klasik ANOVA özellikle grup büyüklükleri eşit değilken ve büyük varyans küçük
grupta bulunduğunda varyans eşitsizliğinden etkilenebilir. Welch ANOVA grup
ortalamalarını kendi tahmin hassasiyetlerine göre ağırlıklandırır ve
serbestlik derecesini ayarlar. Eşit varyans varsayımına gereksinim duymaz.

Welch yaklaşımı şu durumlarda güçlü bir seçenektir:

\begin{itemize}
  \item grup standart sapmaları belirgin farklıysa,
  \item grup büyüklükleri dengesizse,
  \item varyans farklılığı araştırma bağlamında gerçekçi görünüyorsa.
\end{itemize}

Levene veya Brown--Forsythe testi varyans farklılığına ilişkin bilgi verir;
fakat yöntem kararı tek bir eşik üzerinden mekanikleştirilmemelidir. Welch
ANOVA, varyanslar eşit olduğunda da çoğunlukla küçük bir güç kaybıyla çalışır.

\begin{mistakebox}
Varyansların farklı olması ``ANOVA yapılamaz'' anlamına gelmez. Klasik ANOVA
yerine Welch ANOVA, ikili takiplerde ise eşit varyans gerektirmeyen
Games--Howell karşılaştırması kullanılabilir.
\end{mistakebox}

\section{Aykırı değerler ve duyarlılık}
\index{aykırı değer}
\index{duyarlılık analizi}

ANOVA kareli sapmalara dayandığı için uç değerlerden etkilenebilir. Tek bir
puan grup ortalamasını, grup içi kareler toplamını ve dolayısıyla $F$ oranını
değiştirebilir. Aykırı değer adayı görüldüğünde:

\begin{enumerate}
  \item veri girişi ve puanlama kontrol edilir,
  \item gözlemin araştırma bağlamındaki gerçekliği incelenir,
  \item grup grafikleri ve artıklar değerlendirilir,
  \item gerekirse gözlem dahil ve hariç duyarlılık analizi yapılır,
  \item sağlam veya sıralamaya dayalı alternatifler düşünülür.
\end{enumerate}

Gerçek gözlemi yalnız genel testin $p$-değerini değiştirdiği için silmek
seçici analizdir. Ana karar ve bütün alternatif analizler gerekçeleriyle
raporlanmalıdır.

\section{Genel testten sonra hangi gruplar farklı?}
\index{planlı karşılaştırma}
\index{post-hoc karşılaştırma}
\index{Tukey HSD}
\index{Games--Howell testi}

Anlamlı genel $F$ testi farkın yerini göstermez. İki tür takip yaklaşımı
vardır:

\begin{description}
  \item[Planlı karşılaştırma] Kurama dayalı belirli karşılaştırma veri görülmeden önce tanımlanır.
  \item[Post-hoc karşılaştırma] Bütün veya çok sayıdaki ikili fark, aile düzeyindeki hata denetlenerek araştırılır.
\end{description}

Genel test ile takip karşılaştırmasının aynı varsayım ailesinde olması gerekir.
Şekil~\ref{fig:anova-takip-akisi}, ``anlamlı ANOVA'' sonucundan sonra otomatik
olarak bütün ikilileri sınamak yerine araştırma planını ve varyans yapısını
koruyan yolu gösterir.

\begin{figure}[htbp]
\centering
\begin{tikzpicture}[
  font=\small,
  decision/.style={diamond, draw=PrismBlue!80!black, fill=PrismBlue!10,
    align=center, aspect=2.3, inner sep=1.2pt, text width=3.15cm},
  result/.style={draw=PrismBlue!70!black, rounded corners=2pt,
    fill=PrismLight, align=center, minimum height=9mm, text width=3.15cm},
  arrow/.style={-{Stealth[length=2mm]}, thick, draw=PrismBlue!80!black},
  lab/.style={font=\footnotesize, fill=white, inner sep=1pt},
  node distance=7mm and 12mm
]
\node[decision] (omnibus) {Genel test araştırma eşiğini geçti mi?};
\node[result, right=of omnibus] (stop) {İkililer arasında anlamlı olanı arama; tahminleri ve belirsizliği raporla};
\node[decision, below=of omnibus] (planned) {Karşılaştırmalar önceden planlandı mı?};
\node[result, right=of planned] (contrast) {Planlı kontrastlar ve çokluk stratejisi};
\node[decision, below=of planned] (family) {Genel test hangi aileden?};
\node[result] at (-3,-7.3) (tukey) {Klasik ANOVA\\Tukey HSD};
\node[result] at (0,-7.3) (games) {Welch ANOVA\\Games--Howell};
\node[result] at (3,-7.3) (dunn) {Kruskal--Wallis\\Dunn--Holm};
\draw[arrow] (omnibus) -- node[lab, above] {hayır} (stop);
\draw[arrow] (omnibus) -- node[lab, right] {evet} (planned);
\draw[arrow] (planned) -- node[lab, above] {evet} (contrast);
\draw[arrow] (planned) -- node[lab, right] {hayır} (family);
\draw[arrow] (family) -- (tukey);
\draw[arrow] (family) -- (games);
\draw[arrow] (family) -- (dunn);
\end{tikzpicture}
\caption{Genel testten uygun takip karşılaştırmasına geçiş.}
\label{fig:anova-takip-akisi}
\end{figure}

Planlı karşılaştırma, örneğin ``grup danışmanlığı diğer iki düşük yoğunluklu
yaklaşımın ortalamasından farklı mı?'' sorusunu doğrudan sınayabilir. Böyle bir
soru üç ikili karşılaştırmadan daha hedefli olabilir.

Tukey HSD bütün ikili ortalama karşılaştırmalarında aile düzeyindeki Tip I
hata oranını denetler ve klasik ANOVA'nın varyans varsayımıyla uyumludur.
Varyanslar eşit değilse Games--Howell yöntemi daha uygundur. Bonferroni ve
Holm yöntemleri daha genel karşılaştırma ailelerine uygulanabilir
\cite{tukey1953,maxwell2018}.

\begin{table}[htbp]
\centering
\caption{Yaygın takip yöntemlerinin kullanım alanları.}
\begin{tabular}{@{}p{0.22\textwidth}p{0.30\textwidth}p{0.28\textwidth}@{}}
\toprule
Yöntem & Uygun durum & Temel özellik \\
\midrule
Planlı kontrast & Önceden belirlenmiş kuramsal soru & Hedefli ve yorumlanabilir \\
Tukey HSD & Bütün ikililer, benzer varyanslar & Aile düzeyi hatayı denetler \\
Games--Howell & Bütün ikililer, eşit olmayan varyanslar & Welch mantığıyla uyumludur \\
Holm & Sınırlı ve farklı hipotez ailesi & Bonferroni'den genellikle daha güçlüdür \\
\bottomrule
\end{tabular}
\end{table}

\section{Post-hoc sonuçlarının yorumu}

Üç program verisinde Tukey HSD sonuçları şöyledir:

\begin{table}[htbp]
\centering
\caption{Programlar için Tukey ikili karşılaştırmaları.}
\begingroup\scriptsize
\begin{tabular}{@{}p{0.17\textwidth}p{0.14\textwidth}p{0.22\textwidth}p{0.15\textwidth}p{0.15\textwidth}@{}}
\toprule
Karşılaştırma & Ortalama farkı & Yüzde 95 eşzamanlı GA & Düzeltilmiş $p$ & Sonuç \\
\midrule
A eksi B & 3,83 & $[-0{,}59;8{,}25]$ & 0,100 & Belirgin değil \\
A eksi C & 7,33 & $[2{,}91;11{,}75]$ & 0,001 & Belirgin \\
B eksi C & 3,50 & $[-0{,}92;7{,}92]$ & 0,143 & Belirgin değil \\
\bottomrule
\end{tabular}
\endgroup
\end{table}

Genel test fark bulunduğunu göstermiş; düzeltilmiş ikili sonuçlar yalnız A ile
C arasındaki farkı açık biçimde desteklemiştir. A--B ve B--C farklarının
belirgin olmaması bu grupların eşit olduğunu kanıtlamaz. Aralıklar sıfırı ve
uygulamada önemli olabilecek bazı farkları içerdiği için sonuçlar belirsizdir.

\begin{mistakebox}
``A, B'den anlamlı biçimde farklı; B, C'den farklı değil'' sonuçlarından A ile
C hakkında mantıksal bir çıkarım yapılamaz. Her hedef karşılaştırmanın tahmini,
güven aralığı ve düzeltilmiş $p$-değeri doğrudan incelenmelidir.
\end{mistakebox}

\section{Etki büyüklüğü}
\index{eta-kare@$\eta^2$}
\index{omega-kare@$\omega^2$}
\index{etki büyüklüğü!varyans analizi}

Genel $F$ testi farkın varlığına ilişkin kanıtı, etki büyüklüğü ise grup
üyeliğiyle ilişkili değişkenliğin derecesini özetler. Temel ölçülerden
eta-kare
\[
\eta^2=\frac{SS_A}{SS_T}
\]
biçimindedir. Üç program örneğinde
\[
\eta^2=\frac{322{,}89}{965{,}56}\approx0{,}334.
\]
Örneklemde kaygı puanlarındaki toplam değişkenliğin yaklaşık yüzde 33,4'ü grup
üyeliğiyle ilişkilidir.

Eta-kare özellikle küçük örneklemde evren etkisini yukarı yönlü tahmin
edebilir. Daha az yanlı omega-kare
\[
\omega^2=
\frac{SS_A-(k-1)MS_E}{SS_T+MS_E}
\]
ile hesaplanır. Örnekte
\[
\omega^2\approx0{,}288
\]
bulunur.

\begin{researchbox}
``Varyansın yüzde 33'ü program tarafından açıklanmıştır'' ifadesi gözlemsel
bir tasarımda nedensel olarak fazla güçlü olabilir. Daha güvenli ifade,
``örneklemdeki değişkenliğin yüzde 33'ü program gruplarıyla ilişkilidir''
biçimindedir. Açıklama ile nedensellik araştırma tasarımına bağlıdır.
\end{researchbox}

Etki büyüklüğü için sabit küçük--orta--büyük eşikleri bağlamın yerine
geçmemelidir. Ham ortalama farkları ve güven aralıkları, ölçeğin klinik ya da
eğitsel anlamı ve müdahale maliyetiyle birlikte değerlendirilmelidir.

\section{İstatistiksel ve uygulamalı önem}

Genel test anlamlı olsa bile bütün ikili farklar uygulamada önemli olmayabilir.
Örneğin kaygı ölçeğinde hizmet kararını değiştiren en küçük fark 5 puan olsun.
A--C aralığı 2,91 ile 11,75 arasındadır; sıfırı dışlar, fakat 5 puanın hem
altındaki hem üstündeki değerleri içerir. Farkın istatistiksel kanıtı güçlü
olsa da uygulamalı büyüklüğü konusunda belirsizlik sürmektedir.

Tersine, genel test anlamlı değilken geniş güven aralıkları önemli farkları
içerebilir. Bu durumda ``programlar aynıdır'' yerine çalışmanın önemli
farkları dışlayacak hassasiyete sahip olmadığı belirtilmelidir.

\section{Güç ve örneklem hacmi}
\index{test gücü}
\index{örneklem hacmi}

Tek yönlü ANOVA'nın gücü şu etkenlerden etkilenir:

\begin{itemize}
  \item grup ortalamaları arasındaki gerçek farkların büyüklüğü,
  \item grup içi standart sapma,
  \item toplam örneklem ve grup sayısı,
  \item grupların dengeli dağılımı,
  \item anlamlılık düzeyi,
  \item planlı karşılaştırmanın yapısı ve ölçüm güvenirliği.
\end{itemize}

Aynı toplam örneklemde çok sayıda küçük grup oluşturmak her karşılaştırmanın
hassasiyetini düşürebilir. Dengeli grup büyüklükleri genellikle en verimli
tasarımı sağlar. Beklenen kayıp, kümelenme ve çoklu birincil sonuçlar örneklem
planında hesaba katılmalıdır.

Örneklem hacmi yalnız genel $F$ testine göre değil, araştırmanın asıl kararını
oluşturacak planlı karşılaştırma ve uygulamada anlamlı en küçük farka göre
planlanmalıdır.

\section{ANOVA ne zaman yeterli değildir?}
\index{faktöriyel ANOVA}
\index{ANCOVA}
\index{tekrarlı ölçüm}
\index{çok düzeyli model}
\index{etkileşim}

Bu bölümdeki tek yönlü ANOVA bir nicel sonuç ve bir kategorik grup değişkeni
içindir. Daha karmaşık tasarımlar farklı modeller gerektirir:

\begin{itemize}
  \item iki veya daha fazla grup değişkeni: faktöriyel ANOVA ve etkileşim,
  \item ön-test puanını denetleyerek son-test karşılaştırması: ANCOVA veya regresyon,
  \item aynı kişilerin üç veya daha fazla ölçümü: tekrarlı ölçüm veya karma model,
  \item öğrencilerin sınıf ve okul içinde kümelenmesi: çok düzeyli model,
  \item sürekli ve kategorik yordayıcıların birlikte kullanılması: doğrusal regresyon.
\end{itemize}

Özellikle faktöriyel tasarımda ana etkiler, etkileşim varken tek başına
yorumlanmamalıdır. Bir programın etkisi sınıf düzeyine göre değişiyorsa tek bir
ortalama program etkisi bu örüntüyü gizleyebilir.

\begin{etkinlik}
Bir PDR programı, program türü (üç düzey) ve öğrencinin sınıf düzeyi (iki
düzey) bakımından incelenecek olsun. Tek yönlü ANOVA'nın hangi soruları
yanıtlayamayacağını ve olası bir etkileşimin ne anlama gelebileceğini yazın.
\end{etkinlik}

\section{Kruskal--Wallis ve sağlam seçenekler}
\index{Kruskal--Wallis testi}
\index{permütasyon testi}

Dağılımlar ağır çarpık, güçlü aykırı değerli veya ölçüm sıralıysa
Kruskal--Wallis testi düşünülebilir. Bu test ham puanlar yerine sıralara
dayanır. Dağılım biçimleri benzerse konum farkı olarak yorumlanabilir; genel
olarak grupların dağılımlarının aynı olup olmadığını değerlendirir.

Kruskal--Wallis testinin anlamlı olması da hangi grupların farklı olduğunu
göstermez. Dunn gibi düzeltilmiş ikili karşılaştırmalar gerekir. Yöntem
``varsayımsız ANOVA'' değildir; bağımsızlık ve uygun sıralama gibi kendi
koşullarına sahiptir.

Sağlam kırpılmış ortalama yöntemleri, permütasyon testleri ve bootstrap güven
aralıkları da araştırma hedefi ortalama farkı olduğunda yararlı olabilir.
Alternatif seçim istenen sonucu elde etmek için değil, hedef parametre ve veri
yapısına göre yapılmalıdır.

\section{Python ile tek yönlü ANOVA}

Aşağıdaki kod genel testi, ANOVA bileşenlerini, etki büyüklüklerini ve Tukey
karşılaştırmalarını yeniden üretir:

\begin{lstlisting}
import numpy as np
from scipy import stats
from statsmodels.stats.multicomp import pairwise_tukeyhsd

A = np.array([37,49,44,52,41,46,39,50,43,47,42,48])
B = np.array([35,44,40,47,38,43,36,46,39,42,37,45])
C = np.array([31,41,36,44,34,39,32,43,35,40,33,42])
gruplar = [A, B, C]

F, p = stats.f_oneway(*gruplar)
tum = np.concatenate(gruplar)
genel = tum.mean()

ss_grup = sum(len(g) * (g.mean() - genel)**2 for g in gruplar)
ss_hata = sum(((g - g.mean())**2).sum() for g in gruplar)
ss_toplam = ss_grup + ss_hata
df_grup = len(gruplar) - 1
df_hata = len(tum) - len(gruplar)
ms_hata = ss_hata / df_hata

eta2 = ss_grup / ss_toplam
omega2 = ((ss_grup - df_grup * ms_hata) /
          (ss_toplam + ms_hata))

etiket = np.repeat(["A", "B", "C"], [len(A), len(B), len(C)])
tukey = pairwise_tukeyhsd(tum, etiket)

print("Grup ozetleri:", [(g.mean(), g.std(ddof=1)) for g in gruplar])
print("ANOVA:", F, p)
print("Eta-kare ve omega-kare:", eta2, omega2)
print(tukey)
\end{lstlisting}

Kod, klasik ANOVA'yı uygular. Varyanslar belirgin farklıysa Welch ANOVA ve
Games--Howell karşılaştırması için uygun yazılım işlevi seçilmelidir. Sonuçlar
üretilmeden önce grup kodları, eksik gözlemler ve grafikler kontrol edilmelidir.

\section{SPSS ile ANOVA uygulaması}

\begin{enumerate}
  \item \textbf{Analyze $>$ Compare Means $>$ One-Way ANOVA} yolunu izleyin.
  \item Nicel sonucu \textbf{Dependent List}, grup değişkenini \textbf{Factor} alanına taşıyın.
  \item \textbf{Options} bölümünde betimsel istatistikler, varyans homojenliği ve Welch testini seçin.
  \item \textbf{Post Hoc} bölümünde benzer varyanslar için Tukey, eşit olmayan varyanslar için Games--Howell seçeneğini değerlendirin.
  \item Önce grup kodları ve örneklem sayılarını, sonra grafik ve betimsel istatistikleri inceleyin.
  \item Veri yapısına uygun genel testin ardından etki büyüklüğünü ve düzeltilmiş ikili sonuçları raporlayın.
\end{enumerate}

SPSS sürümüne ve kullanılan menüye göre etki büyüklüğü doğrudan verilmeyebilir;
kareler toplamından eta-kare hesaplanabilir. Omega-kare ayrıca formülle
hesaplanmalı veya uygun yazılım yordamından alınmalıdır.

\begin{outputguidebox}
Çıktıyı şu sırayla okuyun: (1) grup tanımları, $n$, ortalama ve standart
sapmalar, (2) kutu grafikleri ve aykırı değerler, (3) varyans yapısı, (4)
klasik veya Welch genel testi, (5) etki büyüklüğü, (6) yalnız gerekiyorsa
uygun düzeltilmiş karşılaştırmalar, (7) fark tahminleri ve eşzamanlı güven
aralıkları. Yalnız ANOVA tablosundaki ``Sig.'' sütununu raporlamayın.
\end{outputguidebox}

\section{Bilimsel raporlama örneği}

``Son-test kaygı puanları program gruplarına göre farklılaşmıştır,
$F(2,33)=8{,}29$, $p=0{,}001$, $\eta^2=0{,}334$,
$\omega^2=0{,}288$. Tukey HSD karşılaştırması A programının ortalamasının C
programından 7,33 puan yüksek olduğunu göstermiştir (yüzde 95 eşzamanlı GA
$[2{,}91;11{,}75]$, düzeltilmiş $p=0{,}001$). A--B ve B--C aralıkları sıfırı
içermiştir. Programlara rastgele atama yapılmadığından sonuçlar nedensel etki
olarak yorumlanmamıştır.''

Bu rapor genel testi, etki büyüklüğünü, farkın yerini, çoklu karşılaştırma
düzeltmesini ve tasarım sınırını birlikte gösterir. Grup ortalamaları ve
standart sapmaları ayrıca tablo veya metinde sunulmalıdır.

\section{Bir ANOVA araştırmasını değerlendirirken}

\begin{enumerate}
  \item Sonuç değişkeni ve grup düzeyleri açıkça tanımlanmış mı?
  \item Gruplar bağımsız mı, eşleştirilmiş mi, yoksa kümelenmiş mi?
  \item Her grubun örneklem büyüklüğü, ortalaması, standart sapması ve grafiği verilmiş mi?
  \item Klasik veya Welch ANOVA seçimi varyans yapısıyla uyumlu mu?
  \item Genel testten sonra karşılaştırmalar önceden planlanmış veya çoklu test için düzeltilmiş mi?
  \item Etki büyüklüğü ve güven aralıkları $p$-değeriyle birlikte sunulmuş mu?
  \item Anlamlı olmayan ikili sonuçlar eşitlik kanıtı gibi yorumlanmış mı?
  \item Araştırma tasarımının izin verdiğinden güçlü nedensel sonuç çıkarılmış mı?
\end{enumerate}

\section{Literatür veri setiyle IBM SPSS laboratuvarı}
\index{Student Performance veri seti!ANOVA}
\index{SPSS!tek yönlü ANOVA}
\index{çalışma süresi!ANOVA uygulaması}
\index{eta-kare!SPSS uygulaması}
\index{omega-kare!SPSS uygulaması}

Bu uygulamada \emph{Student Performance} matematik dosyasındaki yıl sonu
notu \texttt{G3}, öğrencilerin haftalık çalışma süresi gruplarına göre
karşılaştırılacaktır \cite{cortez2008data,cortezsilva2008}. Veri sözlüğünde
\texttt{studytime} dört sıralı kategoriyle kodlanmıştır:

\begin{enumerate}
  \item haftada 2 saatten az,
  \item haftada 2--5 saat,
  \item haftada 5--10 saat,
  \item haftada 10 saatten fazla.
\end{enumerate}

Araştırma sorusu ``ortalama G3 puanı bu dört çalışma süresi grubunda aynı
mıdır?'' biçimindedir. Hedef parametreler dört grubun evren ortalamalarıdır:
\[
H_0:\mu_1=\mu_2=\mu_3=\mu_4,
\qquad
H_1:\text{en az bir grup ortalaması farklıdır}.
\]

\begin{researchbox}
Çalışma süresi öğrencinin kendi bildirimine dayalı sıralı bir kategoridir ve
öğrenciler bu gruplara rastgele atanmamıştır. Ortalama fark bulunsa bile bu
sonuç daha uzun çalışmanın tek başına daha yüksek nota neden olduğunu
kanıtlamaz. Önceki başarı, motivasyon, ders desteği ve okul gibi değişkenler
hem çalışma süresiyle hem notla ilişkili olabilir.
\end{researchbox}

\subsection{Araştırma görevleri}

\begin{enumerate}
  \item Dört grubun kodlarını, örneklem sayılarını, ortalamalarını ve grafiklerini denetleyin.
  \item Klasik tek yönlü ANOVA için bağımsızlık, aykırı değer ve varyans yapısını değerlendirin.
  \item Klasik $F$ testiyle birlikte Welch sağlam testini de isteyin.
  \item Eta-kare ve omega-kare etki büyüklüklerini ANOVA tablosundan hesaplayın.
  \item Genel test anlamlı değilse gerekçesiz post-hoc sonuç aramaktan kaçının.
  \item Sonucu çalışma süresi gruplarının gözlemsel niteliğini belirterek raporlayın.
\end{enumerate}

\subsection{SPSS syntax}

Önce grup etiketleri tanımlanır ve her gruptaki dağılım kutu grafikleriyle
incelenir. \texttt{ONEWAY} komutu betimsel istatistikleri, Levene testini,
klasik ANOVA tablosunu ve Welch sağlam testini birlikte üretir.

\begin{lstlisting}[language={}]
DATASET ACTIVATE StudentMath.

VALUE LABELS studytime
 1 'Haftada 2 saatten az'
 2 'Haftada 2-5 saat'
 3 'Haftada 5-10 saat'
 4 'Haftada 10 saatten fazla'.
VARIABLE LEVEL studytime (ORDINAL).

EXAMINE VARIABLES=G3 BY studytime
 /PLOT BOXPLOT
 /COMPARE GROUPS
 /STATISTICS DESCRIPTIVES
 /CINTERVAL 95
 /MISSING LISTWISE
 /NOTOTAL.

ONEWAY G3 BY studytime
 /STATISTICS DESCRIPTIVES HOMOGENEITY WELCH
 /PLOT MEANS
 /MISSING ANALYSIS.
\end{lstlisting}

Genel test ve veri koşulları ikili karşılaştırmayı gerektirirse eşit
varyanslarda Tukey, eşit varyans varsayımı uygun değilse Games--Howell
istenebilir. Bu veri örneğinde genel test anlamlı olmadığı için bütün ikilileri
yalnız düşük bir $p$-değeri bulmak amacıyla taramak ana analiz planına
eklenmemiştir. Öğretim amacıyla komut yapısı görülmek istenirse ayrı bir
duyarlılık çalışmasında

\begin{lstlisting}[language={}]
ONEWAY G3 BY studytime
 /STATISTICS DESCRIPTIVES HOMOGENEITY WELCH
 /POSTHOC=TUKEY GH ALPHA(.05)
 /MISSING ANALYSIS.
\end{lstlisting}

kullanılabilir; yorumda genel test, varsayım seçimi ve çoklu karşılaştırma
düzeltmesi birlikte belirtilmelidir.

Menüyle çalışmak için \textbf{Analyze $>$ Compare Means $>$ One-Way ANOVA}
yolunda G3 \textbf{Dependent List}, \texttt{studytime} ise \textbf{Factor}
alanına taşınır. \textbf{Options} penceresinde betimsel istatistikler,
varyans homojenliği ve Welch; \textbf{Plots} altında grup ortalamaları seçilir.
Kutu grafikleri için ayrıca \textbf{Analyze $>$ Descriptive Statistics $>$
Explore} kullanılabilir.

\subsection{Grup özetleri}

\begin{table}[htbp]
\centering
\caption{Haftalık çalışma süresine göre G3 grup büyüklükleri ve ortalamaları.}
\begin{tabular}{@{}p{0.44\textwidth}rr@{}}
\toprule
Çalışma süresi grubu & $n$ & Ortalama G3 \\
\midrule
2 saatten az & 105 & 10,048 \\
2--5 saat & 198 & 10,172 \\
5--10 saat & 65 & 11,400 \\
10 saatten fazla & 27 & 11,259 \\
\midrule
Toplam & 395 & 10,415 \\
\bottomrule
\end{tabular}
\end{table}

Örneklem ortalamaları çalışma süresi üçüncü düzeye kadar yükselmekte, dördüncü
düzeyde ise 11,26 civarında kalmaktadır. Ancak dördüncü grupta yalnız 27
öğrenci vardır. Ortalama çizgisinin görünümü tek başına evren ortalamalarının
farklı olduğu sonucunu vermez; grup içi değişkenlik ve örneklem büyüklükleri
genel testte birlikte değerlendirilir.

\subsection{ANOVA çıktısının erişilebilir özeti}

\begin{table}[htbp]
\centering
\caption{Çalışma süresi grupları için klasik tek yönlü ANOVA tablosu.}
\small
\begin{tabular}{@{}lrrrrr@{}}
\toprule
Kaynak & Kareler toplamı & $sd$ & Ortalama kare & $F$ & $p$ \\
\midrule
Gruplar arası & 108,200 & 3 & 36,067 & 1,728 & 0,161 \\
Grup içi & 8161,709 & 391 & 20,874 & & \\
Toplam & 8269,909 & 394 & & & \\
\bottomrule
\end{tabular}
\end{table}

Etki büyüklükleri ANOVA bileşenlerinden
\[
\eta^2=\frac{SS_{\mathrm{grup}}}{SS_{\mathrm{toplam}}}
=\frac{108{,}200}{8269{,}909}\approx0{,}013
\]
ve
\[
\omega^2=
\frac{SS_{\mathrm{grup}}-(k-1)MS_{\mathrm{hata}}}
{SS_{\mathrm{toplam}}+MS_{\mathrm{hata}}}
\approx0{,}005
\]
olarak hesaplanır. SPSS sürümü etki büyüklüğünü doğrudan vermiyorsa aşağıdaki
syntax hesaplamayı belgeler:

\begin{lstlisting}[language={}]
COMPUTE ss_grup=108.200.
COMPUTE ss_toplam=8269.909.
COMPUTE ms_hata=20.874.
COMPUTE k=4.
COMPUTE eta_kare=ss_grup/ss_toplam.
COMPUTE omega_kare=(ss_grup-(k-1)*ms_hata)
 /(ss_toplam+ms_hata).
FORMATS eta_kare omega_kare (F6.3).
LIST eta_kare omega_kare /CASES=FROM 1 TO 1.
\end{lstlisting}

\subsection{Çıktının analizi}

Klasik ANOVA sonucu $F(3,391)=1{,}728$ ve $p=0{,}161$'dir. Dolayısıyla yüzde
5 anlamlılık düzeyinde dört evren ortalamasının eşit olduğu ortak hipoteze
karşı yeterli istatistiksel kanıt bulunmamıştır. Bu sonuç bütün grup
ortalamalarının kesin olarak aynı olduğunu kanıtlamaz; gözlenen farkların grup
içi değişkenliğe göre yeterince belirgin olmadığını gösterir.

Eta-kare yaklaşık 0,013 olduğundan örneklemdeki toplam G3 değişkenliğinin
yaklaşık yüzde 1,3'ü çalışma süresi grupları arasındaki farklarla ilişkilidir.
Yanlılık düzeltmeli omega-kare yaklaşık 0,005'tir. Her iki değer de bu
örneklemde grup ayrımının küçük olduğunu gösterir. Etki büyüklüğü küçük diye
çalışma alışkanlıklarının eğitimsel olarak önemsiz olduğu sonucu çıkarılmaz;
ölçümün geniş kategorilere dayanması ve gözlemsel tasarım da dikkate
alınmalıdır.

Levene ve Welch tabloları varsayım incelemesinin parçalarıdır. Levene sonucu
mekanik bir geçiş anahtarı gibi kullanılmamalı; grup standart sapmaları, kutu
grafikleri, örneklem büyüklüklerindeki dengesizlik ve aykırı değerler birlikte
incelenmelidir. Welch sonucu klasik testten farklı bir karara götürürse,
eşit varyans varsayımına daha az dayanan sonuç öncelikle raporlanmalı ve
farklılığın nedeni açıklanmalıdır.

Genel test anlamlı olmadığı için ana analizde ``hangi ikili anlamlı çıkacak?''
diye altı ayrı karşılaştırma taranmamıştır. Önceden gerekçelendirilmiş belirli
bir planlı karşılaştırma varsa bu, genel testten ayrı bir araştırma hipotezi
olarak açıkça tanımlanabilir; sonuç sonradan seçilmiş gibi sunulmamalıdır.

\begin{mistakebox}
Üçüncü grubun örneklem ortalamasının birinci gruptan yaklaşık 1,35 puan yüksek
olması, iki evren ortalamasının farklı olduğunu tek başına göstermez. Benzer
biçimde $p=0{,}161$ de grupların eşit olduğunu kanıtlamaz. Tahminlerin
belirsizliği, etki büyüklüğü ve çalışma süresinin ölçülme biçimi birlikte
yorumlanmalıdır.
\end{mistakebox}

\begin{outputguidebox}
Çıktıyı şu sırayla okuyun: (1) grup kodları ve geçerli $n$, (2) ortalama,
standart sapma ve güven aralıkları, (3) kutu grafikleri, (4) varyans yapısı ve
Levene testi, (5) klasik ve Welch genel testleri, (6) eta-kare ve omega-kare,
(7) yalnız araştırma planı gerektiriyorsa uygun düzeltilmiş karşılaştırmalar,
(8) gözlemsel tasarımın nedensel yorum sınırı. Anlamlı olmayan genel testten
sonra gerekçesiz ikili sonuç aramayın.
\end{outputguidebox}

\subsection{Örnek bulgu paragrafı}

``Yıl sonu matematik notu ortalamaları haftalık çalışma süresi gruplarında
sırasıyla 10,05 ($n=105$), 10,17 ($n=198$), 11,40 ($n=65$) ve 11,26
($n=27$) olarak bulunmuştur. Klasik tek yönlü ANOVA, grup ortalamalarının
farklı olduğuna ilişkin yeterli kanıt göstermemiştir,
$F(3,391)=1{,}73$, $p=0{,}161$, $\eta^2=0{,}013$,
$\omega^2=0{,}005$. Genel test anlamlı olmadığından gerekçesiz post-hoc
karşılaştırmalar ana analize eklenmemiştir. Çalışma süresi gözlemsel ve geniş
kategorilerle ölçüldüğünden sonuç nedensel etki olarak yorumlanmamıştır.''

\section{R ile çözümlü uygulama}
\label{sec:r-anova}

\textbf{Soru.} Bölümdeki üç grubun $n=(10,10,10)$,
$\bar x=(70,75,82)$ ve $s=(6,7,5)$ özetlerinden klasik ANOVA'yı,
etki büyüklüklerini ve Tukey eşzamanlı fark aralıklarını hesaplayın.
Özetleri ham veri gibi \texttt{aov} komutuna vermeyin.

\begin{lstlisting}[style=prismR]
hacim <- c(10, 10, 10)
ortalama <- c(70, 75, 82)
standart_sapma <- c(6, 7, 5)
grup_sayisi <- length(hacim)
genel <- weighted.mean(ortalama, hacim)
ss_grup <- sum(hacim * (ortalama - genel)^2)
ss_hata <- sum((hacim - 1) * standart_sapma^2)
sd_grup <- grup_sayisi - 1
sd_hata <- sum(hacim) - grup_sayisi
mse <- ss_hata / sd_hata
f_degeri <- (ss_grup / sd_grup) / mse
c(F = f_degeri,
  p = pf(f_degeri, sd_grup, sd_hata, lower.tail = FALSE),
  eta2 = ss_grup / (ss_grup + ss_hata),
  omega2 = (ss_grup - sd_grup * mse) /
    (ss_grup + ss_hata + mse))
ciftler <- combn(seq_along(hacim), 2)
kritik <- qtukey(0.95, nmeans = grup_sayisi, df = sd_hata)
tukey <- t(apply(ciftler, 2, function(cift) {
  fark <- ortalama[cift[1]] - ortalama[cift[2]]
  se <- sqrt(mse / 2 * sum(1 / hacim[cift]))
  c(fark = fark, alt = fark - kritik * se,
    ust = fark + kritik * se,
    p_duzeltilmis = ptukey(abs(fark) / se,
      nmeans = grup_sayisi, df = sd_hata, lower.tail = FALSE))
}))
rownames(tukey) <- c("A-B", "A-C", "B-C")
tukey
\end{lstlisting}

\begin{outputguidebox}
\textbf{Çözüm ve kontrol değerleri.} Genel ortalama 75,6667;
$SS_{\text{grup}}=726{,}6667$, $SS_{\text{hata}}=990$ ve
$MSE=36{,}6667$'dir. $F(2,27)=9{,}9091$, $p=0{,}0005927$,
$\eta^2=0{,}4233$, $\omega^2=0{,}3726$ bulunur. Tukey aralıklarının
yarı genişliği 6,7143'tür. A--B aralığı $[-11{,}7143;1{,}7143]$,
A--C aralığı $[-18{,}7143;-5{,}2857]$, B--C aralığı
$[-13{,}7143;-0{,}2857]$'dir. Son iki aralık sıfırı dışlar.
\end{outputguidebox}

\textbf{Yorum.} Klasik ANOVA ve bu Tukey karşılaştırmaları bağımsız
gözlemler, grup içi yaklaşık normallik ve ortak hata varyansı modeline dayanır.
Genel testin anlamlılığı, üç çiftin hepsinin farklı olduğunu söylemez.
Ham veri bulunduğunda \texttt{aov(puan \textasciitilde{} grup, data = veri)}
ve ardından \texttt{TukeyHSD} kullanılabilir; grup değişkeni faktör olmalıdır.

\subsection*{Eşit olmayan varyanslarda ek çözüm}
Bölümdeki Welch örneğinin özetleriyle aynı düzeltmeyi doğrudan hesaplayın.

\begin{lstlisting}[style=prismR]
hacim <- c(12, 20, 8)
ortalama <- c(70, 76, 84)
varyans <- c(3, 10, 5)^2
grup_sayisi <- length(hacim)
agirlik <- hacim / varyans
pay <- agirlik / sum(agirlik)
merkez <- sum(pay * ortalama)
duzeltme <- sum((1 - pay)^2 / (hacim - 1))
f_welch <- (sum(agirlik * (ortalama - merkez)^2) /
  (grup_sayisi - 1)) /
  (1 + 2 * (grup_sayisi - 2) * duzeltme /
     (grup_sayisi^2 - 1))
sd_payda <- (grup_sayisi^2 - 1) / (3 * duzeltme)
c(F = f_welch, sd_pay = grup_sayisi - 1,
  sd_payda = sd_payda,
  p = pf(f_welch, grup_sayisi - 1, sd_payda,
         lower.tail = FALSE))
\end{lstlisting}

\textbf{Çözüm.} Yaklaşık $F=25{,}3265$, pay serbestlik derecesi 2 ve
payda serbestlik derecesi 18,1621 bulunur; $p=0{,}000005578$'dir. Ham veride
\texttt{oneway.test(puan \textasciitilde{} grup, data = veri, var.equal = FALSE)}
aynı yöntemi uygular. Klasik ANOVA'nın ortak MSE'siyle hesaplanan Tukey
aralıklarını bu eşit olmayan varyans örneğine taşımayın.

\textbf{Pekiştirme ve yanıt.} Klasik örnekte A--B farkı $-5$ olsa da
aralık sıfırı içerir. Bu, iki yöntemin eşdeğer olduğunun kanıtı değildir;
yüzde 5 aile düzeyinde fark gösterilemediğini söyler.
\section{Bölüm sonu çözümlü problemler}

\begin{enumerate}
  \item \textbf{ANOVA tablosunu tamamlama.} $SS_A=240$, $SS_E=600$, $k=4$ ve $N=44$ olsun. Serbestlik derecelerini, ortalama kareleri ve $F$ değerini bulun.

  \textbf{Çözüm:} $sd_A=k-1=3$, $sd_E=N-k=40$'tır. $MS_A=240/3=80$, $MS_E=600/40=15$ ve $F=80/15=5{,}33$ olur. Toplam serbestlik derecesi 43'tür.

  \item \textbf{Eta-kare ve omega-kare.} Önceki problemde $SS_T=840$ olsun. $\eta^2$ ve $\omega^2$ değerlerini bulun.

  \textbf{Çözüm:} $\eta^2=SS_A/SS_T=240/840\approx0{,}286$'dır. $\omega^2=(SS_A-sd_A MS_E)/(SS_T+MS_E)=(240-3\cdot15)/(840+15)=195/855\approx0{,}228$ olur.

  \item \textbf{Takip yöntemi.} Grup büyüklükleri dengesiz ve standart sapmalar belirgin farklıyken genel test anlamlı bulunmuştur. Hangi genel ve ikili yöntemler uygundur?

  \textbf{Çözüm:} Varyans eşitliğine dayanmayan Welch ANOVA uygun genel testtir. Farkın yerini incelemek için Games--Howell ikili karşılaştırmaları kullanılabilir. Genel test tek başına bütün grupların birbirinden farklı olduğunu göstermez.
\end{enumerate}

\bolumkaynaklari{\cite{maxwell2018,kutner2005,cohen1988,tukey1953,cortez2008data,cortezsilva2008}}

\section*{Hatırla--Uygula--Yansıt}

\begin{enumerate}
\renewcommand{\labelenumi}{\textbf{\Alph{enumi}.}}
  \item \textbf{Hatırla:} Tek yönlü ANOVA'nın sıfır ve alternatif hipotezlerini kendi sözlerinizle yazın.
  \item \textbf{Say:} Dört, beş ve altı grubun bütün ikilileri için kaç karşılaştırma gerektiğini hesaplayın.
  \item \textbf{Ayır:} Toplam, gruplar arası ve grup içi değişkenliği bir sınıf örneğiyle açıklayın.
  \item \textbf{Hesapla:} $SS_A=240$, $SS_E=600$, $k=4$, $N=44$ için serbestlik derecelerini, ortalama kareleri ve $F$ değerini bulun.
  \item \textbf{Yorumla:} $F(3,80)=5{,}20$, $p=0{,}003$ sonucunun neyi gösterdiğini ve hangi grupların farklı olduğunu neden göstermediğini açıklayın.
  \item \textbf{Etki:} $SS_A=180$ ve $SS_T=900$ için eta-kareyi hesaplayıp nedensellik iddiası kurmadan yorumlayın.
  \item \textbf{Welch:} Grup büyüklükleri ve standart sapmaları belirgin farklı olduğunda neden Welch ANOVA'nın uygun olabileceğini açıklayın.
  \item \textbf{Post-hoc:} Tukey ile Games--Howell yöntemlerinin kullanım koşullarını karşılaştırın.
  \item \textbf{Aralık:} A--B farkı için eşzamanlı güven aralığı $[-2;7]$ ise ``gruplar eşittir'' sonucunun neden yetersiz olduğunu açıklayın.
  \item \textbf{Eleştir:} Genel test anlamlı çıktıktan sonra düzeltmesiz bütün ikili $t$ testlerini yürütmenin sorununu tartışın.
  \item \textbf{Tasarım:} Aynı öğrencilerin üç zamanda ölçüldüğü bir araştırmada neden bu bölümdeki bağımsız tek yönlü ANOVA'nın uygun olmadığını açıklayın.
  \item \textbf{Duyarlılık:} Tek bir uç puanın $SS_E$, $F$ ve eta-kare üzerindeki olası etkilerini tartışın.
  \item \textbf{Uygula:} Python kodundaki C grubuna iki uç değer ekleyerek genel test, etki büyüklüğü ve Tukey sonuçlarını karşılaştırın.
  \item \textbf{Planla:} Üç program için birincil planlı karşılaştırmayı ve uygulamada önemli en küçük farkı veri görülmeden önce tanımlayın.
  \item \textbf{Raporla:} Grup özetleri, genel test, etki büyüklüğü, düzeltilmiş karşılaştırma, güven aralığı ve tasarım sınırı içeren kısa bir bulgu paragrafı yazın.
\end{enumerate}